\documentclass[12pt]{ltjsarticle}

\usepackage{amsmath,amssymb,amsthm,mathtools}
\usepackage[margin=1in]{geometry}
\usepackage{enumitem}
\usepackage{xurl}
\usepackage[hidelinks]{hyperref}

\theoremstyle{plain}
\newtheorem{theorem}{定理}
\newtheorem{proposition}[theorem]{命題}
\newtheorem{lemma}[theorem]{補題}
\newtheorem{corollary}[theorem]{系}
\theoremstyle{definition}
\newtheorem{remark}[theorem]{注意}
\theoremstyle{plain}
\renewcommand{\proofname}{証明}
\renewcommand{\abstractname}{要旨}

\newcommand{\val}[2]{[#1]_{#2}}
\newcommand{\NpD}{\mathcal{N}_{\mathrm{pD}}}
\newcommand{\DP}{\mathcal{D}_{\mathrm{P}}}
\newcommand{\Pset}{
  \bigl(\tfrac{1}{4}\cdot\mathcal{N}_{pD,\geq 12}\bigr)
}
\newcommand{\seqnum}[1]{\href{https://oeis.org/#1}{\textrm{\underline{#1}}}}

\title{原始Dyck数による加法的表現：\\例外集合と最適位数}
\author{Takayuki Kuriyama}
\date{}

\begin{document}

\maketitle

\begin{abstract}
$\NpD$ を、$0$ と、原始Dyck語を2進数として読んで得られる整数とからなる集合とする。正の偶数 $N$ に対し、$r(N)$ を、和が $N$ となる $\NpD$ の正の元の最小個数とする。このとき、対応する数列 $a(n)=r(2n)$ はOEIS数列 \seqnum{A395858} である。本論文では、6項表示に関する完全な例外集合を決定する。すなわち、$r(46)=8$ を満たすのは $46$ のみであり、10個の整数
$34,44,98,154,198,202,206,838,842,846$ は $r(N)=7$ を満たし、それ以外のすべての正の偶数は $r(N)\leq6$ を満たす。したがって、6項で常に表示できるようになる鋭い偶数閾値は $848$ である。さらに、正確な相対漸近位数も決定する。任意の $k\geq2$ に対して$  r(10\cdot4^{k+1}-6)=6$
が成り立つので、無限に多くの場合に6項が必要であり、漸近位数は正確に $6$ である。証明では、正規部分言語、有限区間証明書、自己相似的な5項区間伝播、および2色Motzkin語による半分化原始族の基数4特徴づけを組み合わせる。さらに、その半分化族が正確に位数 $5$ の漸近加法基底であることも従う。
\end{abstract}

\noindent\textbf{キーワード：} 原始Dyck数、加法的表現、加法基底、漸近位数、例外集合、2色Motzkin語、ディジタル自己相似性、区間充填。

\noindent\textbf{2020 Mathematics Subject Classification:} Primary 11B13;
Secondary 05A15, 68Q45.

\section{序論}\label{sec:introduction}

均衡した括弧列は、開き括弧を $1$、閉じ括弧を $0$ とすることにより、アルファベット $\{1,0\}$ 上の語として書くことができ、したがって2進整数として読むことができる。この方法でどの整数が得られるのか、また、それらがなす集合は加法的にどれほど豊かであるのか。本論文では、この問いの自然な変種を扱う。

\emph{Dyck語}とは、$1$ と $0$ を同数含み、かつ任意の接頭辞において $1$ の個数が $0$ の個数以上である語 $w\in\{1,0\}^*$ のことである。ここで $1$ は開き括弧、$0$ は閉じ括弧を表す。空でないDyck語が\emph{原始的}であるとは、その空でない真の接頭辞のどれも均衡していないことをいう。同値に、対応する格子路が終点において初めて高さ $0$ に戻ることをいう。このような路は、\emph{既約}Dyck路または\emph{分解不能}Dyck路とも呼ばれる。原始Dyck語の言語を $\DP$ と書く。長さ順、同じ長さでは辞書式順に並べると、最初の18語は
\[
\begin{gathered}
10,\ 1100,\ 110100,\ 111000,\ 11010100,\ 11011000,\\
11100100,\ 11101000,\ 11110000,\ 1101010100,\ 1101011000,\\
1101100100,\ 1101101000,\ 1101110000,\ 1110010100,\\
1110011000,\ 1110100100,\ 1110101000
\end{gathered}
\]
である。

基数 $m$ の数字からなる語 $w$ に対し、それが基数 $m$ で表す整数を $\val{w}{m}$ と書く。加法表現のために便宜上 $0$ を加えて、
\[
 \NpD=\{0\}\cup\{\val{w}{2}:w\in\DP\}
\]
と定めるが、これは定理の結果に影響を与えない。最初の18個の正の元は
\[
\begin{gathered}
2,12,52,56,212,216,228,232,240,\\
852,856,868,872,880,916,920,932,936
\end{gathered}
\]
であり、これはOEIS数列 \seqnum{A057547} \cite{OEIS-A057547} である。空でないDyck語はすべて末尾が $0$ であるため、$\NpD$ の正の元はすべて偶数である。したがって、$\NpD$ の元の和で表示できる対象整数は偶数に限られる。

以下の、より鋭い合同式に関する観察を繰り返し用いる。

\begin{lemma}\label{lem:mod4}
$4$ を法として $2$ に合同な正の原始Dyck数は $2$ だけである。それ以外の正の原始Dyck数はすべて $4$ で割り切れる。
\end{lemma}

\begin{proof}
空でないDyck語はすべて末尾が $0$ である。さらに、長さが4以上の原始Dyck語は実際には末尾が $00$ でなければならない。もし末尾が $10$ ならば、最後の2文字の直前で既に路が高さ $0$ に戻ってしまうからである。語 $10$ は $2$ を表し、したがってそれより長い原始語はすべて $4$ の倍数を表す。
\end{proof}

Bell、Lidbetter、Shallitは、\emph{すべての}Dyck語から得られる、より大きな集合である\emph{完全均衡数}（OEIS数列 \seqnum{A014486} \cite{OEIS-A014486}）を研究した。彼らは正規な下方近似とオートマトンを用い、
\[
        8,\ 18,\ 28,\ 38,\ 40,\ 82,\ 166
\]
を除くすべての偶数が高々3個の完全均衡数の和であり、一方で漸近的には2項では不十分であることを証明した \cite[Thm.\ 14]{BLS}。

原始性という制限は、この結果の表面的な変種ではない。空でないDyck語はすべて、原始Dyck語の連結として一意に分解されるが、本論文では各加数そのものが単一の因子でなければならない。したがって、制限のないDyck数による表示から、一般には同じ項数の原始Dyck数表示は得られない。

数字条件で定義された加法基底は、オートマトン理論的手法によっても研究されてきた。Bell、Hare、Shallitは、加法基底となる自動集合を特徴づけ、その位数を決定する有効な手続きを与えた \cite{BHS}。関連する決定手続きにより、2進回文数 \cite{RSS} および2進平方語 \cite{MNRS} についても鋭い加法的結果が得られている。原始Dyck語の言語は正規ではないため、これらの一般的な有限オートマトン手法だけでは本問題を直接解決できない。

本論文の証明は、正規近似だけでは不十分になる構造的な地点で、従来の方法と異なる。正規部分言語が $4$ を法として $0$ の合同類を処理し、$4$ を法として $2$ の合同類は、$4$ で割った後の完全な原始族がもつ自己相似的閉性によって処理する。同じ族は、2色Motzkin語による基数4表示ももつ。その自己相似性は相補的な二つの効果をもつ。すなわち、5項区間を伝播させる一方で、世代間の隙間が4項表示に対する無限個の障害を生み出す。これらの事実から、完全例外集合、最適な一様上界、鋭い最終的閾値、および正確な漸近位数が得られる。

長さ $2n$ のDyck語のうち原始的なものの割合は、$n\to\infty$ のとき $1/4$ に収束する。実際、長さ $2n$ のDyck語は $C_n$ 個あり、原始Dyck語は一意に $1u0$ と書ける。ここで $u$ は長さ $2n-2$ のDyck語である。したがって、長さ $2n$ の原始Dyck語は $C_{n-1}$ 個である。Catalan数の公式
$C_n=\frac1{n+1}\binom{2n}{n}$
から
\footnote{文法 $S\to\varepsilon\mid 1S0S$ により、空でないDyck語は一意に $1u0v$（$u,v$ はDyck語）と分解される。 したがって、その個数 $C_n$ は $C_0=1$ および $C_n=\sum_{k=0}^{n-1}C_kC_{n-1-k}$ を満たす。 一方、$B_n=\frac{1}{n+1}\binom{2n}{n}$ とおけば、二項係数の畳み込み公式から $B_0=1$ および $B_n=\sum_{k=0}^{n-1}B_kB_{n-1-k}$ が成り立つ。 ゆえに帰納法により $C_n=B_n$、すなわち $C_n=\frac{1}{n+1}\binom{2n}{n}$ である。}
、
\[
 \frac{C_{n-1}}{C_n}=\frac{n+1}{2(2n-1)}\longrightarrow\frac14
\]
を得る。しかし、その算術的効果は著しい。十分大きな偶数はすべて、制限のないDyck数なら3項で足りるのに対し、原始族には小さいながら非自明な例外集合が存在する。本論文では、高々6個の原始Dyck数の和でない正の偶数が正確に11個であることを示す。そのうち10個は7項を必要とし、$46$ だけが8項を必要とする。これら11個の例外はすべて $848$ 未満にある。

8項を必要とする現象は完全に局所的である。$48$ 未満の正の原始Dyck数は $2$ と $12$ しかない。一方、最終的な6項上界は二つの自己相似族から生じる。正規部分族が $4$ の倍数を処理し、$4$ で割った完全原始族が残る合同類を処理する。

正の偶数 $N$ に対し、和が $N$ となる正の原始Dyck数の最小個数を $r(N)$ とし、本論文の中心的な整数列を
\[
        a(n)=r(2n),\qquad n\geq1
\]
と定める。この数列はOEISに \seqnum{A395858} \cite{OEIS-A395858} として登録されている。最初の50項は
\[
\begin{gathered}
1,2,3,4,5,1,2,3,4,5,6,2,3,4,5,6,7,3,4,5,6,7,8,4,5,\\
1,2,1,2,3,4,2,3,2,3,4,5,3,4,3,4,5,6,4,5,4,5,6,7,5
\end{gathered}
\]
である。次のPythonプログラムは、ここに示した項を再現し、数列の50,000項を生成して、いくつかの反復パターンを可視化する：
\href{https://github.com/growupkuriyama-hub/lean_cfg_project/blob/3525cfb/LeanCfgProject/PrimDyck/A395858.py}{固定されたリポジトリ・スナップショット \texttt{3525cfb} における \path{A395858.py}}。

\begin{theorem}\label{thm:classification}
$r(N)$ が6を超える場合は、正確に次のとおりである：
\[
\begin{aligned}
 r(46)&=8,\\[4pt]
 r(N)=7\quad&\Longleftrightarrow\quad
 N\in\{34,44,98,154,198,202,206,838,842,846\}.
\end{aligned}
\]
したがって、それ以外のすべての正の偶数は $r(N)\leq6$ を満たす。特に、任意の偶数 $N\geq848$ は高々6個の原始Dyck数の和である。
\end{theorem}

同値に、$a(n)=r(2n)$ で定義されるOEIS数列 \seqnum{A395858} \cite{OEIS-A395858} について、
\[
        \max_{n\geq1}a(n)=8,
        \qquad
        a(n)=8\quad\Longleftrightarrow\quad n=23
\]
であり、
\[
 a(n)=7
 \quad\Longleftrightarrow\quad
 n\in\{17,22,49,77,99,101,103,419,421,423\}.
\]
したがって、主定理は単なる最終的評価ではなく、$a(n)$ が6を超えるすべての項を完全に決定している。

最終的な上界6もまた最良である。

\begin{theorem}\label{thm:asymptotic-order}
\[
 h_{\mathrm{asym}}
 =\min\{h:\text{十分大きなすべての偶数 }N\text{ が }r(N)\leq h\text{ を満たす}\}
\]
とおく。このとき $h_{\mathrm{asym}}=6$ である。より正確には、任意の整数 $k\geq2$ に対して
\[
 N_k=10\cdot4^{k+1}-6
\]
は $r(N_k)=6$ を満たす。同値に、
\[
 a(5\cdot4^{k+1}-3)=6\qquad(k\geq2)
\]
である。
\end{theorem}

したがって、上記の有限例外集合と、6項上界が達成される無限に多くの大きな整数とは両立する。

例外値 $46$ だけでも、一様上界が8より小さくなり得ないことが分かる。

\begin{proposition}\label{prop:46}
整数 $46$ は正確に8個の正の原始Dyck数を必要とする。すなわち、8個の和ではあるが、高々7個の和ではない。
\end{proposition}

\begin{proof}
$48$ 未満の正の原始Dyck数は $2$ と $12$ だけである。実際、長さ2および4の原始語は $10$ と $1100$ である一方、長さ6以上の原始語は $11$ で始まるので、少なくとも $[110000]_2=48$ を表す。したがって、$46$ の任意の表示は
\[
        46=12a+2b,\qquad a,b\geq0
\]
の形である。同値に $6a+b=23$ である。$12\cdot4>46$ なので $a\leq3$ であり、したがって項数は
\[
        a+b=23-5a\geq8
\]
を満たす。等号は
\[
        46=12+12+12+2+2+2+2+2
\]
によって達成される。
\end{proof}

例えば、
\[
 2026=2+240+852+932=2+216+872+936
\]
である。この二つの表示に現れる7個の相異なる整数は、すべて原始Dyck語を2進展開にもつ。

以下で用いる和集合の記法は、加法基底論における標準的なものである \cite[Chap.\ 1]{Nathanson}。非負整数 $k$ に対し、$kX$ を、$X$ の元を重複を許して $k$ 個足して得られる和の集合とし、$0X=\{0\}$ とする。また、
\[
 F_k(X)=\bigcup_{j=0}^{k}jX
\]
とおく。したがって、$F_k(X)$ は $X$ の元を高々 $k$ 個足して得られる和の集合である。スカラー $c$ に対しては、$c\cdot X=\{cx:x\in X\}$ と書き、これは $X$ の拡大を表す。

法4の性質から、中心的な半分化族
\begin{equation}\label{eq:P-def}
 \Pset=\{n/4:v\in\NpD,\ n\geq12\}
\end{equation}
を導入する。このとき
\begin{equation}\label{eq:NpD-decomp}
 \NpD=\{0,2\}\cup4\Pset
\end{equation}
である。$\Pset$ の最初の元は
\[
 3,13,14,53,54,57,58,60,213,214,217,\ldots
\]
である。

\section{\texorpdfstring{$\Pset$}{P} の基数\texorpdfstring{$4$}{4}構造}\label{sec:base4}

4進数の各桁に重みを
\[
  \sigma(3)=1,\qquad \sigma(0)=-1,\qquad
  \sigma(1)=\sigma(2)=0
\]
によって割り当てる。語 $w=c_0\cdots c_{k-1}\in\{0,1,2,3\}^{k}$ が\emph{2色Motzkin語}であるとは、任意の接頭辞の総重みが非負であり、かつ $w$ 全体の総重みが $0$ であることをいう。重み0の二つの数字 $1$ と $2$ は、異なる色とみなす。$w$ を基数4の数として読んだ値を $[w]_4$ と書き、$[\varepsilon]_4=0$ とする。

\begin{lemma}[基数4による特徴づけ]\label{lem:base4}
正整数 $p$ が $\Pset$ に属するための必要十分条件は、その4数展開が $3w$ の形であり、$w$ が2色Motzkin語であることである。
\end{lemma}

\begin{proof}
$p\in\Pset$とする。$4p\in\NpD$かつ$4p\geq12$, つまり$\operatorname{bin}(4p)\geq 1100$である。$w=\operatorname{bin}(4p)=d_{0}\cdots d_{k-1}$, $d_i\in\{0,1\}^2$, $k\geq4$とする。 語$w$を左から読むとき、組$d_i$をよむ直前における2進路の高さ$h(d_0\cdots d_{i-1})$は、累積重み$\sigma(d_0\cdots d_{i-1})$の2倍である。実際、
\[
  11\leftrightarrow3,\quad 10\leftrightarrow2,\quad
  01\leftrightarrow1,\quad 00\leftrightarrow0
\]
だからである。$4p$は$12$以上の原始Dyck数なので $\operatorname{bin}(4p)$は先頭が $11$、末尾が $00$ である。したがって 4進数解釈が$[d_0]_4=3$、$[d_{k-1}]_4=0$ である。最後以外の組$d_{j}$を読む時点$d_0\cdots d_{j-1}$では$h(d_0\cdots d_{j-1})\geq2$ であり、読み終えた最後で $h(d_0\cdots d_{k-1})=0$ となる。

これらの境界条件は、組$b_j$の内部の高さも制御する。型 $01$ の組は、高さを一時的に1だけ下げてから元に戻す。最後でない型 $00$ の組は、終了時の高さが少なくとも $2$ なので、$b_{j-1}$を読む時点での高さは少なくとも $4$ である。また、残る型$11,\, 10$はいずれも $1$ で始まる。したがって、最後の$b_{k-1}$を読む前で高さが正であることは、最後の$b_{k-1}$を読む時点で $h=0$ に戻ることと併せれば、全ての途中地点$j<k$で$h(h_0\cdots h_{j-1})>0$であることと同値である。

$v$ を $4$ で割ることは、末尾の基数4数字 $0$ を削除することに対応する。したがって、$p$ の基数4展開は $3w$ であり、$3w$ の任意の接頭辞の $\sigma$ 重みは少なくとも $1$、また $3w$ 全体の重みは $1$ である。先頭の数字 $3$ を取り除けば、$w$ の任意の接頭辞の重みが非負であり、$w$ 全体の重みが $0$ であることが分かる。$w$ が2色Motzkin語とわかった。

逆に、$p$ の基数4展開が $\operatorname{quad}(p)=3w\in\{0,1,2,3\}^*$ で、$w$ が2色Motzkin語であるとする。末尾に数字 $0$ を付け加え、各基数4数字を対応する2進の2桁組に展開する。得られる2進語$11\operatorname{bin}([w]_4)00=:b_0\cdots b_{k-1}$, $b_i\in\{0,1\}^2$は $1$ と $0$ を同数含み、先述の$\sigma=2h$の関係性から最後のビット$b_{k-1}$より前で$b_{j-1}$を読む時点では累積高さ$h(b_0\cdots b_{j-1})$は$0$ より真に大である。したがって原始Dyck語であり、$4p\in\NpD$、ゆえに $p\in\Pset$ である。
\end{proof}

この特徴づけから、後で用いる自己相似性が得られる。

\begin{corollary}\label{cor:Pclosure}
$p\in\Pset$ ならば、$4p+1$ および $4p+2$ も $\Pset$ に属する。
\end{corollary}

\begin{proof}
$p$ の基数4展開を $\operatorname{quad}(p)=3w\in\{0,1,2,3\}^*$ と書く。ただし $w$ は2色Motzkin語である。このとき、 $\operatorname{quad}(4p+1)=3w1$ ,  $\operatorname{quad}(4p+2)=3w2$ である。重み0のいずれの色$\operatorname{quad}([1]_4)=01$, $\operatorname{quad}([2]_4)=10$を末尾に付け加えても2色Motzkin条件は保たれるので、補題~\ref{lem:base4} から主張が従う。
\end{proof}

次に、各世代における最小値と最大値を決定する。

\begin{lemma}[2色Motzkin値の極値]\label{lem:extremal}
任意の $\ell\geq0$ と、長さ $\ell$ の任意の2色Motzkin語 $w$ に対して、
\begin{equation}\label{eq:motzkin-extrema}
  \frac{4^{\ell}-1}{3}\leq[w]_4\leq4^{\ell}-2^{\ell}
\end{equation}
が成り立つ。両方の境界は達成される。下界は $1^{\ell}$ によって達成され、上界は
\[
  w_{\max,\ell}=
  \begin{cases}
    3^a0^a,&\ell=2a,\\
    3^a2\,0^a,&\ell=2a+1
  \end{cases}
\]
によって達成される。
\end{lemma}

\begin{proof}
$\ell$ に関する帰納法を用いる。$\ell=0$ の場合は明らかである。

\emph{下界。}
$\sigma(0)=-1$のため、最初の数字は $0$ ではあり得ない。$w=c_0\cdots c_{\ell-1}$, $c_j\in\{0,1,2,3\}$とする。$\sigma(1)=\sigma(2)=0$のため、$c_0\in\{1,2\}$ の場合は、接尾辞 $w'=c_1\cdots c_{\ell-1}$ も2色Motzkin語なので、
\[
  [w]_4\geq4^{\ell-1}+[w']_4
  \geq4^{\ell-1}+\frac{4^{s-1}-1}{3}
  =\frac{4^{\ell}-1}{3}
\]
である。また、$\sigma(3)=2$であるので、$c_0=3$ の場合は、
\[
  [w]_4\geq [3\underbrace{00\cdots0}_{\ell-1\text{ 桁}}]_4=3\cdot4^{\ell-1}\geq\frac{4^{\ell}-1}{3}
\]
である。等号は $w=1^{\ell}$ によって達成される。

\emph{上界。}
$c_0\in\{1,2\}$ の場合、$w'$ も2色Motzkin語なので、
\[
  [w]_4
  =[c_0\underbrace{c_1c_2\cdots c_{\ell-1}}_{\ell-1\text{ 桁}}]_4\leq[2\underbrace{c_1c_2\cdots c_{\ell-1}}_{\ell-1\text{ 桁}}]_4\leq2\cdot4^{\ell-1}+\bigl(4^{\ell-1}-2^{\ell-1}\bigr)
  =3\cdot4^{\ell-1}-2^{\ell-1}
  \leq4^{\ell}-2^{\ell}
\]
である。

次に $c_0=3$ とし、接頭辞の重みが初めて $0$ に戻る位置を $r-1$ とする。このとき $1\leq r\leq\ell$、$c_r=0$ であり、
\[
  u=c_1\cdots c_{r-1},\qquad v=c_{r}\cdots c_{\ell-1}
\]
は、それぞれ長さ $r-2$ と $\ell-r$ の2色Motzkin語である。位取り記数法と帰納法の仮定から、
\begin{align*}
  [w]_4
  &=[3\underbrace{c_1c_2\cdots c_{r-1}}_{r-1\text{ 桁}}\underbrace{c_{r}c_{r+1}\cdots c_{\ell-1}}_{\ell-r\text{ 桁}}]_4\\
  &=3\cdot4^{\ell-1}+[c_1c_2\cdots c_{r-1}]_4\,4^{\ell-r+1}+[c_{r}c_{r+1}\cdots c_{\ell-1}]_4\\
  &\leq3\cdot4^{\ell-1}
    +\bigl(4^{r-2}-2^{r-2}\bigr)4^{\ell-r+1}
    +4^{\ell-r}-2^{\ell-r}\\
  &=4^{\ell}-2^{2\ell-r}+2^{2\ell-2r}-2^{\ell-r}\\
    &=4^{\ell}-(2^{2\ell-r}-2^{2\ell-2r})-2^{\ell-r}
\end{align*}
を得る。残る確認は
\[
  2^{2\ell-r}-2^{2\ell-2r}
  \geq2^{\ell}-2^{\ell-r}
\]
である。因数分解すると、これは
\[
  2^{2\ell-2r}(2^r-1)\geq2^{\ell-r}(2^r-1)
\]
となり、$r\leq\ell$ から従う。これで上界が証明された。

最後に、表示した語 $w_{\max,\ell}$ は2色Motzkin語である。等比数列の和を計算すれば、
\[
  [w_{\max,\ell}]_4=4^{\ell}-2^{\ell}
\]
を得るので、上界は達成される。
\end{proof}

$\Pset$ の元の基数4展開がちょうど $d$ 桁であるとき、その元は\emph{世代 $d$} に属するという。補題~\ref{lem:base4} により、世代 $d$ の元の値は $3\cdot4^{d-1}+[w]_4$ である。ここで $w$ の長さは $d-1$ である。

\begin{corollary}[世代境界]\label{cor:gen}
任意の $d\geq1$ に対し、$\Pset$ の世代 $d$ の任意の元は
\begin{equation}\label{eq:generation-interval}
  [g_d,U_d]
  =\left[
  3\cdot4^{d-1}+\frac{4^{d-1}-1}{3},
  4^d-2^{d-1}
  \right]
\end{equation}
に属し、両端点はとも達成される。特に、
\begin{enumerate}[label=\textup{(\alph*)}]
\item $(U_d,3\cdot4^d)$ には $\Pset$ の元が存在しない。
\item 任意の $k\geq0$ に対して、
\begin{equation}\label{eq:three-g}
  3g_{k+1}=10\cdot4^k-1
\end{equation}
が成り立つ。
\end{enumerate}
\end{corollary}

\begin{proof}
区間の主張は補題~\ref{lem:base4} と補題~\ref{lem:extremal} から従う。端点が達成されることは、補題~\ref{lem:extremal} の極値語から従う。(a) は、次の世代が $g_{d+1}>3\cdot4^d$ から始まることによる。最後に、
\[
  3g_{k+1}
  =3\left(3\cdot4^k+\frac{4^k-1}{3}\right)
  =10\cdot4^k-1
\]
である。
\end{proof}

\section{正規族と6項区間充填}\label{sec:regular}

第~\ref{sec:base4} 節で記述した基数4の族の内部には、原始Dyck語を貫く単純な正規パターンがある。例えば、$868$ の2進展開は
\[
 \operatorname{bin}(868)=1101100100=11\,01\,10\,01\,00
 \in 11(01\mid10)^*00
\]
と分解される。2進数字を2桁ずつ組にすると、最初のブロック $11$ は基数4の数字 $3$ に、最後のブロック $00$ は $0$ に、各中間ブロック $01$ または $10$ は $1$ または $2$ に変わる。これは、そのような基数4整数を固定個数足した和が長い区間を埋める可能性を示唆する。以下の混合和集合補題によって、この考えを正確に定式化する。

\[
 \mathcal E=\{0\}\cup
 \bigl\{\val{3\varepsilon_1\cdots\varepsilon_k}{4}:
        k\geq0,\ \varepsilon_i\in\{1,2\}\bigr\}
\]
とおく。$k=0$ のとき、文字列 $\varepsilon_1\cdots\varepsilon_k$ は空なので、対応する元は $\val{3}{4}=3$ である。$4$ を掛けることは、末尾に基数4の $0$ を付け加えることに対応する。各基数4数字を対応する2進2桁に置き換えると、
\[
\begin{aligned}
 3_{(4)}&\longleftrightarrow11_{(2)},&
 1_{(4)}&\longleftrightarrow01_{(2)},\\
 2_{(4)}&\longleftrightarrow10_{(2)},&
 0_{(4)}&\longleftrightarrow00_{(2)}
\end{aligned}
\]
である。したがって、$4\cdot\mathcal E$ の正の各元は、2進展開が $11(01\mid10)^*00$ に属する。最初の $11$ の後で路は高さ $2$ にあり、各中間ブロック $01$ または $10$ は路を高さ $2$ に戻し、最後の $00$ が初めて高さ $0$ に戻す。ゆえに、$4\cdot\mathcal E$ の正の各元は原始Dyck数であり、$0\in\NpD$ なので、
\begin{equation}\label{eq:4E}
        4\cdot\mathcal E\subseteq\NpD
\end{equation}
である。

核心は、$\mathcal E$ の6個の元だけで $25$ 以上のすべての整数を覆うことである。

\begin{proposition}\label{prop:E}
任意の整数 $n\geq25$ は $6\mathcal E$ に属する。
\end{proposition}

これを有限桁の区間論法によって証明する。まず、以下の補助集合の役割を説明しておく。$\mathcal E_m$ の $m$ 桁の元を先頭桁と下位桁に分けると、下位 $m-1$ 桁は、各桁が $\{0,1\}$ に属する基数4の補正をなす。集合 $\Delta_m$ は、これらの補正を正確に記録する。6項和のうちいくつかの加数が新しい先頭桁を用いると、残る下位桁の問題は、$\Delta_{m-1}$ と $\mathcal E_{m-1}$ のコピーの混合和となる。混合和集合 $S_{t,m}$ は、これらすべての可能性を同時に追跡するために設計されている。

$u_0=0$ とおき、$m\geq1$ に対して
\begin{equation}\label{eq:umdef}
 u_m=[\underbrace{11\cdots1}_{m\text{ 桁}}]_4
    =1+4+\cdots+4^{m-1}=\frac{4^m-1}{3}
\end{equation}
と定める。このとき、
\begin{equation}\label{eq:urecursion}
\begin{aligned}
 4^{m-1}&=3u_{m-1}+1 &&(m\geq1),\\
 u_{m-1}&=4u_{m-2}+1 &&(m\geq2)
\end{aligned}
\end{equation}
である。補正集合を
\[
 \Delta_m=\Bigl\{[\delta_{m-1}\cdots\delta_0]_4:\delta_j\in\{0,1\}\Bigr\}
    =\Bigl\{\textstyle\sum_{j=0}^{m-1}\delta_j4^j:\delta_j\in\{0,1\}\Bigr\}
\]
と定める。このとき $\max \Delta_m=u_m$ である。$\mathcal E_m$ を、$0$ と、基数4で高々 $m$ 桁の元とからなる $\mathcal E$ の有限部分集合とする。例えば、
\[
\begin{aligned}
 \Delta_2&=\{[00]_4,[01]_4,[10]_4,[11]_4\}=\{0,1,4,5\},\\
 \mathcal E_2&=\{0,[3]_4,[31]_4,[32]_4\}=\{0,3,13,14\}
\end{aligned}
\]
である。$0\leq t\leq6$ に対して
\begin{equation}\label{eq:S-t-m}
 S_{t,m}=t\Delta_m+(6-t)\mathcal E_m
\end{equation}
とおく。$\mathcal E_m$ の最大元は
\[
 \max\mathcal E_m=[3\underbrace{22\cdots2}_{m-1\text{ 桁}}]_4
                =3\cdot4^{m-1}+2u_{m-1}
\]
であり、$m\geq2$ のとき、式~\eqref{eq:urecursion} から
\begin{equation}\label{eq:maxErecursion}
 \max\mathcal E_{m-1}=11u_{m-2}+3=\frac{11u_{m-1}+1}{4}
\end{equation}
を得る。$\max \Delta_m=u_m$ なので、
\begin{equation}\label{eq:maxStm}
 \max S_{t,m}=tu_m+(6-t)\max\mathcal E_m
\end{equation}
である。さらに $\max\mathcal E_m>u_m$ なので、$\max S_{t,m}$ は $t$ とともに減少する。以下、$[a,b]$ は整数区間 $\{a,a+1,\ldots,b\}$ を表す。

\begin{lemma}\label{lem:mixed}
任意の $m\geq2$ および任意の $2\leq t\leq6$ に対して、
\begin{equation}\label{eq:mixedtge}
 S_{t,m}=[0,\max S_{t,m}]
\end{equation}
が成り立つ。一方、
\begin{equation}\label{eq:mixedtone}
 S_{1,m}=[0,\max S_{1,m}]\setminus\{2\}
\end{equation}
である。さらに、
\begin{equation}\label{eq:mixedEminterval}
 [25,\ell_m]\subseteq6\mathcal E_m,
 \qquad \ell_m=\frac{33}{8}\,4^m-4
\end{equation}
が成り立つ。
\end{lemma}

\begin{proof}
まず基底の場合 $m=2$ において、\eqref{eq:mixedtge} と \eqref{eq:mixedtone} を証明する。先に見たように、
\[
 \Delta_2=\{0,1,4,5\},\qquad \mathcal E_2=\{0,3,13,14\}
\]
である。直接計算すると、$2\Delta_2=\{0,1,2,4,5,6,8,9,10\}$ であり、さらに $\Delta_2$ を一つ加えると残りのすべての隙間が埋まる。したがって $3\Delta_2=[0,15]$ であり、ゆえに
\[
 t\Delta_2=[0,5t]\qquad(3\leq t\leq6)
\]
である。一方、
\[
\begin{aligned}
2\mathcal E_2
={}&\{0,3,6,13,14,16,17,26,27,28\},\\[1ex]
3\mathcal E_2
={}&\{0,3,6,9\}\cup[13,14]\cup[16,17]\cup[19,20]\\
&\qquad\cup[26,31]\cup[39,42],\\[1ex]
4\mathcal E_2
={}&\{0,3,6,9\}\cup[12,14]\cup[16,17]\cup[19,20]\\
&\qquad\cup[22,23]\cup[26,34]\cup[39,45]\cup[52,56],\\[1ex]
5\mathcal E_2
={}&\{0,3,6,9\}\cup[12,17]\cup[19,20]\cup[22,23]\\
&\qquad\cup[25,37]\cup[39,48]\cup[52,59]\cup[65,70]
\end{aligned}
\]
である。$5\mathcal E_2$ に $\Delta_2$ を加え、$4\mathcal E_2$ に $2\Delta_2$ を加えると、
\[
 S_{1,2}=\{0,1\}\cup[3,75],\qquad S_{2,2}=[0,66]
\]
を得る。同様に、
\[
\begin{aligned}
S_{3,2}&=3\Delta_2+3\mathcal E_2=[0,15]+3\mathcal E_2=[0,57],\\[1ex]
S_{4,2}&=4\Delta_2+2\mathcal E_2=[0,20]+2\mathcal E_2=[0,48],\\[1ex]
S_{5,2}&=5\Delta_2+\mathcal E_2=[0,25]+\mathcal E_2=[0,39],\\[1ex]
S_{6,2}&=6\Delta_2=[0,30]
\end{aligned}
\]
である。まとめると、
\[
\begin{array}{c|c}
t&S_{t,2}\\ \hline
1&[0,75]\setminus\{2\}\\
2&[0,66]\\
3&[0,57]\\
4&[0,48]\\
5&[0,39]\\
6&[0,30]
\end{array}
\]
となり、$m=2$ における \eqref{eq:mixedtge} と \eqref{eq:mixedtone} が証明された。\eqref{eq:mixedEminterval} については、$[25,\ell_2]=[25,62]\subseteq6\mathcal E_2$ を確認すればよい。区間 $[26,28]$、$[39,42]$、$[52,56]$ は、それぞれ $\{13,14\}$ の元を2個、3個、4個足して得られるので、
\[
\begin{aligned}
25&=13+3+3+3+3+0,\\
[26,40]&=\{26,27,28\}+3\cdot\{0,1,2,3,4\},\\
[39,51]&=[39,42]+3\cdot\{0,1,2,3\},\\
[52,62]&=[52,56]+3\cdot\{0,1,2\}
\end{aligned}
\]
である。これで基底の場合が完了する。下端点は鋭い。実際、$25$ 未満の $\mathcal E$ の正の元は $3,13,14$ のみであり、これらを高々6個足しても $24$ にはならないので、$24\notin6\mathcal E$ である。

次に $m\geq3$ とし、レベル $m-1$ で \eqref{eq:mixedtge}--\eqref{eq:mixedEminterval} が成り立つと仮定する。この段階の記法を軽くするため、
\[
 u=u_{m-1},\qquad
 q=q_m=[3\underbrace{11\cdots1}_{m-1\text{ 桁}}]_4
      =3\cdot4^{m-1}+u,\qquad
 M=\max\mathcal E_{m-1}
\]
とおく。このとき
\[
 4^{m-1}=3u+1,\qquad
 M=\frac{11u+1}{4}
\]
である。また、
\begin{equation}\label{eq:qmdef}
 q=10u+3
\end{equation}
である。以下、繰り返し
\begin{equation}\label{eq:qoverlap}
 2\cdot4^{m-1}+4u=q-1
\end{equation}
を用いる。実際、両辺はともに $10u+2$ である。基数4の先頭桁に従って分解すると、
\[
 \Delta_m=\Delta_{m-1}\cup(4^{m-1}+\Delta_{m-1}),\qquad
 \mathcal E_m=\mathcal E_{m-1}\cup(q+\Delta_{m-1})
\]
を得る。実際、$\Delta_m$ の新しい $m$ 桁の元は先頭桁が $1$ であり、一方、$\mathcal E_m$ の新しい $m$ 桁の元は先頭桁が $3$ で、その後に $\Delta_{m-1}$ から来る下位 $m-1$ 桁の補正が続く。$t$ 個の $\Delta_m$ のうち $s$ 個が平行移動された部分を用い、$6-t$ 個の $\mathcal E_m$ のうち $v$ 個が新しい $m$ 桁部分を用いると、\eqref{eq:S-t-m} の和集合は
\begin{equation}\label{eq:decomp}
 S_{t,m}=\bigcup_{v=0}^{6-t}\ \bigcup_{s=0}^{t}
 \bigl(s\cdot4^{m-1}+vq+S_{t+v,m-1}\bigr)
\end{equation}
と分解される。したがって、$v$ を固定すると、パラメータ $s$ は $4^{m-1}$ の倍数だけ平行移動した $t+1$ 個の集合を生じ、$v$ を増やすと次の大きなブロックへ移る：
\[
 \underbrace{[vq,\ \cdots]}_{v\text{-ブロック}}
 \qquad
 \underbrace{[(v+1)q,\ \cdots]}_{(v+1)\text{-ブロック}}.
\]
以下では、隣接する平行移動どうし、および隣接する $v$-ブロックどうしが重なることを確認する。

まず $t\geq2$ の場合を扱う。帰納法の仮定 \eqref{eq:mixedtge} により、\eqref{eq:decomp} に現れるすべての $v$ に対して、$S_{t+v,m-1}=[0,\max S_{t+v,m-1}]$ である。$v$ を固定すると、連続する $s$ の値はこの区間を $4^{m-1}$ ずつ平行移動する。$\max S_{j,m-1}$ は $j$ とともに減少するので、
\[
 \max S_{t+v,m-1}\geq\max S_{6,m-1}=6u
     \geq3u=4^{m-1}-1
\]
である。したがって、これら $t+1$ 個の平行移動は重なり、
\[
 [vq,\,vq+t\cdot4^{m-1}+\max S_{t+v,m-1}]
\]
を形成する。連続する $v$ に対応するブロックどうしも重なる。次のブロックが存在するときは $t+v\leq5$ であり、$t\geq2$ なので、
\[
\begin{aligned}
 t\cdot4^{m-1}+\max S_{t+v,m-1}
 &\geq2\cdot4^{m-1}+\max S_{5,m-1}\\
 &=2\cdot4^{m-1}+5u+M\\
 &=(2\cdot4^{m-1}+4u)+(u+M)\\
 &\geq q-1
\end{aligned}
\]
である。最後の不等式には \eqref{eq:qoverlap} を用いた。ゆえに \eqref{eq:decomp} の和集合は、$0$ から最大元 $\max S_{t,m}$ までの全整数区間である。これでレベル $m$ における \eqref{eq:mixedtge} が証明された。

$t=1$ の場合、最初のブロックには欠けた整数 $2$ が一つだけ残る。各固定された $v\geq1$ に対し、二つの $s$-平行移動に現れるのは $S_{1+v,m-1}$ であり、$1+v\geq2$ である。それらは
\[
 \max S_{1+v,m-1}\geq\max S_{6,m-1}=6u\geq4^{m-1}-1
\]
により重なる。したがって、後続の各 $v$-ブロックは区間である。$\max S_{j,m-1}$ は $j$ とともに減少するので、連続する $v$-ブロックの最後の組だけを確認すれば十分である。実際、
\[
\begin{aligned}
 4^{m-1}+\max S_{5,m-1}-(q-1)
 &=M-2u-1\\
 &=\frac{11u+1}{4}-2u-1
  =\frac{3u-3}{4}\geq0
\end{aligned}
\]
なので、連続する $v$-ブロックは重なる。最後に、
\[
\begin{aligned}
 \max S_{1,m-1}-(4^{m-1}+2)
 &=u+5M-4^{m-1}-2\\
 &=\frac{47u-7}{4}\geq0
\end{aligned}
\]
である。したがって、最初の $v$-ブロック内で、平行移動されていない長い区間は $4^{m-1}+2$ まで達し、次の $s$-平行移動と隙間なくつながる。ゆえに、
\[
 S_{1,m}=[0,\max S_{1,m}]\setminus\{2\}
\]
を得る。これでレベル $m$ における \eqref{eq:mixedtone} が証明された。

残るのは、$6\mathcal E_m$ の内部で区間を伝播させることである。\eqref{eq:decomp} で $t=0$ とすると、
\[
 6\mathcal E_m=\bigcup_{v=0}^{6}\bigl(vq+S_{v,m-1}\bigr)
\]
である。$v=0$ のブロックは $S_{0,m-1}=6\mathcal E_{m-1}$ であり、帰納法の仮定により $[25,\ell_{m-1}]$ を含む。$m\geq3$ なので $u\geq5$ であり、
\[
 \ell_{m-1}-(q+2)=\frac{19u-39}{8}\geq0
\]
である。$v=1$ のブロック $q+S_{1,m-1}$ は、帰納法の仮定により $\{q,q+1\}\cup[q+3,q+\max S_{1,m-1}]$ を含む。直前の区間は $q+2$ まで達しているので、両者は隙間なくつながる。$j=1,2,3$ に対し、$v=j$ のブロックと $v=j+1$ のブロックをつなぐには、
\[
 \max S_{j,m-1}\geq q-1
\]
であれば十分である。$\max S_{j,m-1}$ は $j$ とともに減少するので、最も厳しい場合は $j=3$ である。したがって、
\[
\begin{aligned}
 \max S_{3,m-1}-(q-1)
 &=3u+3M-(10u+2)\\
 &=\frac{5u-5}{4}\geq0
\end{aligned}
\]
を確認すればよい。ゆえに $\max S_{1,m-1},\ \max S_{2,m-1},\ \max S_{3,m-1}\geq q-1$ であり、ブロック $v=1,2,3,4$ は順に連結して、
\[
 4q+\max S_{4,m-1}=\frac{33}{8}\,4^m-4=\ell_m
\]
までを覆う。残るブロック $v=5,6$ は隙間の後から始まる可能性があるが、ここで伝播させる区間には不要である。これでレベル $m$ における \eqref{eq:mixedtge}--\eqref{eq:mixedEminterval} が証明され、帰納法が完了する。
\end{proof}

$\ell_m\to\infty$ かつ $\mathcal E=\bigcup_{m\geq1}\mathcal E_m$ なので、命題~\ref{prop:E} は直ちに従う。

上の命題は、$4$ を法として $0$ の合同類を処理する。すなわち、$100$ 以上の任意の $4$ の倍数は、高々6個の原始Dyck数の和である。

\section{完全な半分化族に対する5項区間}
\label{sec:five-interval}

正規族 $\mathcal E$ は $4$ の倍数を扱うには十分であるが、もう一方の合同類に必要な5項区間を、それ単独では与えない。ここでは、\eqref{eq:P-def} で定義した完全な族 $\Pset$ と、系~\ref{cor:Pclosure} の自己相似的閉性を用いる。

以下、$\Pset$ に対する5項区間を確立する。有限な基底部分には、二つの部分集合
\begin{equation}\label{eq:L-def}
 L=\{3,13,14,53,54,57,58,60\}\subseteq\Pset
\end{equation}
および
\begin{equation}\label{eq:H-def}
\begin{split}
 H=\{&213,214,217,218,220,229,230,233,234,236,\\
     &241,242,244,248\}\subseteq\Pset
\end{split}
\end{equation}
を用いる。次の表は、対応するすべての原始Dyck数を示している。各欄において、2進語 $[4h]_2$ は原始的であり、したがって $h\in\Pset$ である。
\[
\begin{array}{c|c@{\qquad}c|c}
\multicolumn{4}{c}{4L}\\
h&[4h]_2&h&[4h]_2\\ \hline
3&1100&53&11010100\\
13&110100&54&11011000\\
14&111000&57&11100100\\
58&11101000&60&11110000
\end{array}
\]
\[
\begin{array}{c|c@{\qquad}c|c}
\multicolumn{4}{c}{4H}\\
h&[4h]_2&h&[4h]_2\\ \hline
213&1101010100&233&1110100100\\
214&1101011000&234&1110101000\\
217&1101100100&236&1110110000\\
218&1101101000&241&1111000100\\
220&1101110000&242&1111001000\\
229&1110010100&244&1111010000\\
230&1110011000&248&1111100000
\end{array}
\]

有限計算は、Pythonプログラム \path{verify_certificates_for_paper.py} によって再現できる。このプログラムは、\href{https://github.com/growupkuriyama-hub/lean_cfg_project/blob/3525cfb/LeanCfgProject/PrimDyck/verify_certificates_for_paper.py}{コミット \texttt{3525cfb} に固定されたリポジトリ版}として利用できる。保存された版は、そのコミットにおけるリポジトリのPython CI検査を通過している（Python CI run \#1）。同じプログラムは、実行コマンドと期待される出力を記したREADMEとともに、補助ファイルとしても添付される。

補題~\ref{lem:finitefive} の5つの主張は包含証明書である。すなわち、プログラムは、表示された各区間が対応する和集合に含まれることを検査するのであり、その和集合が区間と等しいとは主張しない。一方、補題~\ref{lem:finiteexception} については、プログラムは $[0,211]$ に切り詰めた関連集合を正確に計算し、表示された区間和集合との等式を検査する。

\begin{lemma}[有限区間証明書]\label{lem:finitefive}
次の包含関係が成り立つ：
\begin{equation}\label{eq:finite-interval-certificate}
\begin{array}{c|c}
\text{和集合}&\text{含まれる整数区間}\\ \hline
5L &[215,254]\\
H+4L &[245,486]\\
2H+3L &[439,674]\\
3H+2L &[645,862]\\
4H+L &[855,1046].
\end{array}
\end{equation}
\end{lemma}

\begin{proof}
\eqref{eq:finite-interval-certificate} の最初の包含、すなわち $[215,254]\subseteq5L$ を示す。\eqref{eq:L-def} で定義した集合 $L$ を
\[
 A=\{53,54,57,58,60\},\qquad B=\{3,13,14\}
\]
と分割し、$L=A\sqcup B$ とする。順次加えると、
\[
 2A=[106,108]\cup[110,118]\cup\{120\}
\]
および
\[
 3A=[159,178]\cup\{180\},\qquad
 4A=[212,238]\cup\{240\}
\]
を得る。したがって $4A+B=[215,254]$ であり、\eqref{eq:finite-interval-certificate} の最初の包含が従う。

\eqref{eq:finite-interval-certificate} の第二の包含 $[245,486]\subseteq H+4L$ について、$L$ を繰り返し加えると、
\begin{equation}\label{eq:4L-intervals}
\begin{split}
4L\supseteq{}&[32,34]\cup[42,45]\cup[52,56]\cup[82,102]\\
 &{}\cup[116,148]\cup[162,194]\cup[212,238]
\end{split}
\end{equation}
を得る。\eqref{eq:4L-intervals} の右辺にある7区間のうち、最初の3つの短い区間に \eqref{eq:H-def} の集合 $H$ を加えると、
\[
\begin{aligned}
H+[32,34]&\supseteq[245,254]\cup[261,270]
 \cup[273,278]\cup[280,282],\\
H+[42,45]&\supseteq[255,265]\cup[271,281]\cup[283,293],\\
H+[52,56]&\supseteq[265,276]\cup[281,304]
\end{aligned}
\]
を得る。これらの和集合は $[245,304]$ を含む。$\min H=213$、$\max H=248$ であり、$H$ の連続する元の差は高々 $9$ なので、$I=[a,b]$ が $b-a\geq8$ を満たせば、
\[
 H+I=[213+a,248+b]
\]
である。これを \eqref{eq:4L-intervals} の最後の4区間に適用すると、それぞれ
\[
 [295,350],\quad[329,396],\quad[375,442],\quad[425,486]
\]
を得る。これらの区間は互いに重なるので、\eqref{eq:finite-interval-certificate} の第二の包含 $[245,486]\subseteq H+4L$ の証明が完了する。

\eqref{eq:finite-interval-certificate} の第三の包含 $[439,674]\subseteq2H+3L$ については、次の含まれる区間で十分である：
\[
 2H\supseteq[426,428]\cup[430,438]\cup[446,478]\cup\{496\},
\]
および
\[
\begin{split}
3L\supseteq{}&\{9\}\cup[19,20]\cup[39,42]\cup[73,77]
 \cup[109,111]\\
 &{}\cup[113,134]\cup[159,178]\cup\{180\}.
\end{split}
\]
$2H$ の最初の二つの成分から、
\[
 [430,438]+9=[439,447],\quad
 [426,428]+[19,20]=[445,448],
\]
および
\[
 [430,438]+[19,20]=[449,458]
\]
を得る。$[446,478]$ に、順に
\[
\begin{gathered}
 \{9\},\ [39,42],\ [73,77],\ [109,111],\\
 [113,134],\ [159,178],\ \{180\}
\end{gathered}
\]
を加えると、$[455,658]$ を覆う互いに重なる区間が得られる。最後に、
\[
 496+[159,178]=[655,674]
\]
である。以上の区間を合わせると、第三の包含が従う。

\eqref{eq:finite-interval-certificate} の第四の包含 $[645,862]\subseteq3H+2L$ については、
\[
 3H\supseteq[639,734]\cup\{744\},\qquad
 2L\supseteq\{6\}\cup[70,74]\cup[110,118]\cup\{120\}
\]
を用いる。四つの区間和
\[
 [645,740],\quad[709,808],\quad[759,854],\quad[854,862]
\]
は $[645,862]$ を覆う。

\eqref{eq:finite-interval-certificate} の第五の包含 $[855,1046]\subseteq4H+L$ については、
\[
 4H\supseteq[852,982]\cup[984,986],
 \qquad \{3,57,58,60\}\subseteq L
\]
を用いる。このとき、区間
\[
 [855,985],\quad[912,1042],\quad[1041,1044],\quad[1044,1046]
\]
が $[855,1046]$ を覆う。
\end{proof}

補題~\ref{lem:finitefive} の5区間は互いに重なるので、
\[
 [215,1046]\subseteq5\Pset
\]
を得る。

\begin{proposition}\label{prop:fiveP}
任意の整数 $n\geq215$ は、$\Pset$ の元を正確に5個足した和である。したがって、
\[
 [212,\infty)\subseteq F_5(\Pset)
\]
である。下端点は鋭い。すなわち、$209,210,211$ のいずれも $F_5(\Pset)$ に属さない。
\end{proposition}

\begin{proof}
補題~\ref{lem:finitefive} から、より強い包含 $[215,1046]\subseteq5\Pset$ が得られる。$n$ に関する強帰納法の基底として必要なのは、初期区間 $[215,867]$ だけである。$n\geq868$ とし、$[215,n-1]$ のすべての整数が $5\Pset$ に属すると仮定する。$\rho\equiv n\pmod4$ を満たす一意な $\rho\in\{5,6,7,8\}$ を選ぶ。このとき、
\[
 M=\frac{n-\rho}{4}\geq\frac{868-8}{4}=215
\]
である。$M<n$ なので、帰納法の仮定により
\[
 M=q_1+\cdots+q_5,\qquad q_i\in\Pset
\]
と書ける。$5\leq\rho\leq8$ なので、
\[
 \eta_1+\cdots+\eta_5=\rho
\]
を満たす $\eta_i\in\{1,2\}$ を選ぶことができる。したがって、
\[
 n=4M+\rho=\sum_{i=1}^{5}(4q_i+\eta_i)
\]
である。系~\ref{cor:Pclosure} により、各加数 $4q_i+\eta_i$ は $\Pset$ に属する。ゆえに $n\in5\Pset$ であり、帰納法が完了する。

最後に、
\[
\begin{aligned}
212&=53+53+53+53,\\
213&=53+53+53+54,\\
214&=53+53+54+54
\end{aligned}
\]
なので、$[212,\infty)\subseteq F_5(\Pset)$ である。次に、$212$ 未満の $\Pset$ の元が $L$ の8元だけであることを確認する。長さ4、6、8の原始Dyck語を $4$ で割ると、それぞれ
\[
 \{3\},\qquad \{13,14\},\qquad \{53,54,57,58,60\}
\]
が得られる。長さ10の原始Dyck語はすべて $1u0$ の形であり、ここで $u$ は長さ8のDyck語である。辞書式順で最小の語は $1101010100$ で、その値は $852$ である。$4$ で割ると $213$ を得る。より長い原始語の値はさらに大きい。したがって、\eqref{eq:L-def} の $L$ に対して
\[
 \Pset\cap[0,211]=L
\]
である。

\[
 D=\{53,54,57,58,60\}\subseteq L
\]
とおく。$L$ の元を高々5個足した和で、$D$ の元を高々3個しか用いないものは、最大でも
\[
 3\cdot60+2\cdot14=208
\]
である。一方、$D$ の元を少なくとも4個用いる和は、最小でも
\[
 4\cdot53=212
\]
である。したがって $209,210,211\notin F_5(\Pset)$ である。
\end{proof}

\section{完全例外集合}\label{sec:exceptions}

まず、$4$ を法として $2$ の合同類における6項表示可能性を、$\Pset$ に関する有限和条件へ変換する。

\begin{lemma}\label{lem:sixcriterion}
$N=4m+2$ とする。このとき、$N$ が高々6個の正の原始Dyck数の和であるための必要十分条件は、
\[
 m\in F_5(\Pset)
 \cup\bigl(1+F_3(\Pset)\bigr)
 \cup\bigl(2+F_1(\Pset)\bigr)
\]
である。
\end{lemma}

\begin{proof}
補題~\ref{lem:mod4} により、$2$ 以外の原始Dyck加数はすべて $4$ の倍数である。したがって、$N\equiv2\pmod4$ の表示には、奇数個の $2$ が含まれる。項数が高々6なら、その個数は $1,3,5$ のいずれかである。

$2$ が1個の場合、残りの和を $4$ で割ると $m\in F_5(\Pset)$ を得る。$2$ が3個の場合は $m-1\in F_3(\Pset)$、5個の場合は $m-2\in F_1(\Pset)$ を得る。これで必要性が示された。逆に、これら三つの包含のいずれからも、対応する $\Pset$ の加数を $4$ 倍し、それぞれ $2$ を1個、3個、5個付け加えることにより $N$ の表示が得られる。したがって条件は十分でもある。
\end{proof}

\begin{lemma}[小さい4の倍数]\label{lem:smallmultiples}
$100$ 未満の正の $4$ の倍数のうち、高々6個の正の原始Dyck数で表示できないものは $44$ だけである。さらに、
\[
 44=12+12+12+2+2+2+2
\]
である。
\end{lemma}

\begin{proof}
$100$ 未満の正の原始Dyck数は正確に
\[
 2,\ 12,\ 52,\ 56
\]
である。まず $N<52$ とし、$N=4x$ と書く。$12$ のコピーと、対 $2+2$ の和は
\[
 x=3a+c
\]
の形であり、$a+2c$ 個の加数を用いる。条件 $a+2c\leq6$ のもとで得られる $x<13$ の正の値は
\[
 1,2,3,4,5,6,7,8,9,10,12
\]
である。唯一欠ける値は $x=11$ で、これは $N=44$ に対応する。

次に $52\leq N<100$ とする。まず $52=4\cdot13$ または $56=4\cdot14$ のいずれかを用いる。残り高々5項で、$12$ および対 $2+2$ から得られる剰余商は
\[
 R=\{3a+c:a,c\geq0,\ a+2c\leq5\}
   =\{0,1,2,3,4,5,6,7,9,10,12,15\}
\]
である。これら12個の整数を直接調べると、
\[
 [13,24]\subseteq(13+R)\cup(14+R)
\]
である。したがって、$52$ から $96$ までのすべての $4$ の倍数は、求める表示をもつ。上に表示した $44$ の7項表示により証明が完了する。
\end{proof}

\begin{lemma}[有限例外集合証明書]\label{lem:finiteexception}
\[
 L=\{3,13,14,53,54,57,58,60\}
\]
とする。このとき、
\[
\begin{aligned}
F_5(L)\cap[0,211]
={}&\{0,3,6,9\}\cup[12,17]\cup[19,20]\cup[22,23]\\
 &{}\cup[25,37]\cup[39,48]\cup[52,208]
\end{aligned}
\]
である。さらに、
\[
\begin{aligned}
[0,211]\cap\Bigl(F_5(L)&\cup(1+F_3(L))\cup(2+F_1(L))\Bigr)\\
={}&[0,7]\cup[9,10]\cup[12,23]\cup[25,37]\\
 &{}\cup[39,48]\cup[52,208]
\end{aligned}
\]
である。したがって、$[0,211]$ における補集合は
\[
 \{8,11,24,38,49,50,51,209,210,211\}
\]
である。
\end{lemma}

\begin{proof}
$0L=\{0\}$ から始め、漸化式 $jL=(j-1)L+L$ を用いる。$L$ の8元を順次加えると、$211$ で切り詰めた次の表が得られる：
\[
\begin{array}{c|l}
j & F_j(L)\cap[0,211]\\
\hline
1 & \{0,3\}\cup[13,14]\cup[53,54]\cup[57,58]\cup\{60\}\\
2 & \{0,3,6\}\cup[13,14]\cup[16,17]\cup[26,28]\cup[53,54]\cup[56,58]\\
  & {}\cup[60,61]\cup\{63\}\cup[66,68]\cup[70,74]\cup[106,108]\cup[110,118]\cup\{120\}\\
3 & \{0,3,6,9\}\cup[13,14]\cup[16,17]\cup[19,20]\cup[26,31]\cup[39,42]\cup[53,54]\\
  & {}\cup[56,61]\cup[63,64]\cup[66,77]\cup[79,88]\cup[106,134]\cup[159,178]\cup\{180\}\\
4 & \{0,3,6,9\}\cup[12,14]\cup[16,17]\cup[19,20]\cup[22,23]\cup[26,34]\cup[39,45]\\
  & {}\cup[52,64]\cup[66,102]\cup[106,148]\cup[159,194]\\
5 & \{0,3,6,9\}\cup[12,17]\cup[19,20]\cup[22,23]\cup[25,37]\cup[39,48]\cup[52,208].
\end{array}
\]
また、
\[
\begin{aligned}
(1+F_3(L))\cap[0,211]
={}&\{1,4,7,10\}\cup[14,15]\cup[17,18]\cup[20,21]\\
 &{}\cup[27,32]\cup[40,43]\cup[54,55]\cup[57,62]\\
 &{}\cup[64,65]\cup[67,78]\cup[80,89]\cup[107,135]\\
 &{}\cup[160,179]\cup\{181\}
\end{aligned}
\]
であり、
\[
 (2+F_1(L))\cap[0,211]
 =\{2,5\}\cup[15,16]\cup[55,56]\cup[59,60]\cup\{62\}
\]
である。表示した区間リストの和集合を取ると第二の式が得られ、そこから欠ける整数は最後に表示した集合と正確に一致する。表の各行は、直前の行を $L$ の8元だけそれぞれ平行移動し、その和集合を直前の行自身と合わせることによって得られる。
\end{proof}

\begin{proof}[定理~\ref{thm:classification} の証明]
まず $N\equiv0\pmod4$ とする。$N\geq100$ なら $N=4n$ と書く。このとき $n\geq25$ なので、命題~\ref{prop:E} と \eqref{eq:4E} により、$N$ は高々6個の原始Dyck数で表示できる。補題~\ref{lem:smallmultiples} は残る正の $4$ の倍数を処理し、$r(44)=7$ を示す。

次に $N\equiv2\pmod4$ とし、$N=4m+2$ と書く。$N\geq850$ なら $m\geq212$ なので、命題~\ref{prop:fiveP} により $m\in F_5(\Pset)$ である。補題~\ref{lem:sixcriterion} により、$N$ は高々6個の原始Dyck数の和である。前段落と合わせると、任意の偶数 $N\geq848$ がそのような表示をもつことが既に分かる。

残るのは、$N\leq846$ かつ $N\equiv2\pmod4$ の整数を分類することである。この場合 $0\leq m\leq211$ であり、命題~\ref{prop:fiveP} から、$m$ 以下の $\Pset$ の元は
\[
 L=\{3,13,14,53,54,57,58,60\}
\]
に属するものだけである。補題~\ref{lem:finiteexception} により、判定条件で覆われない $m$ の値は
\[
 \{8,11,24,38,49,50,51,209,210,211\}
\]
である。補題~\ref{lem:sixcriterion} から、対応する $N=4m+2$ の値は正確に
\[
 34,46,98,154,198,202,206,838,842,846
\]
である。したがって、この10個と $44$ を合わせた11個が、高々6個の原始Dyck数の和でない正の偶数のすべてである。

$46$ 以外はすべて7項表示をもつ：
\[
\begin{aligned}
34&=12+12+2+2+2+2+2,\\
44&=12+12+12+2+2+2+2,\\
98&=56+12+12+12+2+2+2,\\
154&=52+52+12+12+12+12+2,\\
198&=56+52+52+12+12+12+2,\\
202&=56+56+52+12+12+12+2,\\
206&=56+56+56+12+12+12+2,\\
838&=240+240+240+52+52+12+2,\\
842&=240+240+240+56+52+12+2,\\
846&=240+240+240+56+56+12+2
\end{aligned}
\]
である。命題~\ref{prop:46} により、$46$ は正確に8項を必要とする。これですべての主張が証明された。
\end{proof}

\begin{corollary}\label{cor:threshold}
$848$ は、任意の偶数 $N\geq T$ が高々6個の原始Dyck数の和となるような最小の偶数 $T$ である。
\end{corollary}

\begin{proof}
定理~\ref{thm:classification} により、$848$ 以上のすべての偶数では6項で十分であり、一方 $846$ は7項を必要とする。
\end{proof}

\begin{corollary}\label{cor:eight}
任意の非負偶数は $\NpD$ の高々8個の元の和であり、$46$ は $\NpD$ の高々7個の元の和でない唯一の偶数である。
\end{corollary}

\begin{proof}
正の整数については定理~\ref{thm:classification} から直ちに従い、整数 $0$ は空和である。
\end{proof}

$A\subseteq2\mathbb N_0$ に対し、$2\mathbb N_0$ の任意の元が $A$ の高々 $h$ 個の元の和であるとき、$A$ を $2\mathbb N_0$ に対する\emph{位数高々 $h$ の相対加法基底}と呼ぶ。この性質をもつ最小の整数が $h$ であるとき、その相対位数は\emph{正確に $h$} であるという。同様に、$2\mathbb N_0$ の十分大きな任意の元が $A$ の高々 $h$ 個の元の和であるとき、$A$ を\emph{位数高々 $h$ の相対漸近加法基底}と呼ぶ。

定理~\ref{thm:classification} は6項表示に関する完全例外集合を与える。系~\ref{cor:eight} は $\NpD$ が $2\mathbb N_0$ に対する正確に位数8の相対加法基底であることを示し、系~\ref{cor:threshold} は鋭い6項閾値を与える。次節では、相対漸近位数が正確に6であることを証明する。

\section{反復する隙間と漸近的最適性}\label{sec:asymptotic}

$4$ を法として $2$ に合同な整数の表示は、例外的な加数 $2$ の個数に従って分かれる。5項表示に対する判定条件は次のとおりである。

\begin{lemma}[5項判定条件]\label{lem:fivecrit}
$m\geq3$ とし、$N=4m+2$ とする。このとき $r(N)\leq5$ であるための必要十分条件は、
\begin{equation}\label{eq:fivecrit}
  m\in F_4(\Pset)\cup\bigl(1+F_2(\Pset)\bigr)
\end{equation}
である。
\end{lemma}

\begin{proof}
補題~\ref{lem:mod4} により、$N$ の表示は $2$ をある個数 $s\geq0$ 個用い、残りは $4\Pset$ の元を用いる。法 $4$ で考えると $2s\equiv2\pmod4$ なので、$s$ は奇数である。高々5項を用いるなら、$s\in\{1,3,5\}$ である。残りの加数を $4p_i$ と書き、$N=2s+\sum_i4p_i$ から $2$ を引いて $4$ で割ると、
\[
  m=\frac{s-1}{2}+\sum_i p_i
\]
を得る。$s=1$ の場合は $m\in F_4(\Pset)$、$s=3$ の場合は $m\in1+F_2(\Pset)$ となる。$s=5$ なら $m=2$ となるが、これは $m\geq3$ に反する。これで必要性が示された。逆に、\eqref{eq:fivecrit} の最初の集合に属していれば $2$ を1個用いる表示が得られ、第二の集合に属していれば $2$ を3個用いる表示が得られる。
\end{proof}

次の補題は、下界を与える機構を取り出したものである。

\begin{lemma}[反復障害]\label{lem:obstruction}
任意の整数 $k\geq2$ に対し、
\[
  x_k=10\cdot4^k-2
\]
とおく。このとき、
\begin{equation}\label{eq:obstruction}
  x_k\notin F_4(\Pset)
  \qquad\text{かつ}\qquad
  x_k-1\notin F_2(\Pset)
\end{equation}
である。
\end{lemma}

\begin{proof}
$k\geq2$ を固定し、$x=x_k$ と略記する。世代 $k+2$ の最小元は
\[
  g_{k+2}>3\cdot4^{k+1}=12\cdot4^k>x
\]
である。したがって、$x$ 以下の $\Pset$ の元はすべて世代 $k+1$ 以下に属する。系~\ref{cor:gen} により、そのような元は、世代 $k$ 以下であるため
\[
  U_k=4^k-2^{k-1}
\]
以下である\emph{小さい}元か、世代 $k+1$ に属し、$g_{k+1}$ 以上 $U_{k+1}=4^{k+1}-2^k$ 以下である\emph{大きい}元のいずれかである。

まず $x-1=10\cdot4^k-3$ を考える。これは $U_{k+1}$ より大きく、$g_{k+2}$ より小さいので、それ自身は $\Pset$ に属さない。さらに、許される二つの元の和は最大でも
\[
  2U_{k+1}=8\cdot4^k-2^{k+1}<10\cdot4^k-3=x-1
\]
である。したがって $x-1\notin F_2(\Pset)$ である。

次に、背理法のため
\[
  x=q_1+\cdots+q_t,\qquad t\leq4,\qquad q_i\in\Pset
\]
と仮定する。大きい加数の個数を $\lambda$ とする。$\lambda\leq2$ なら、残る位置を可能な限り大きい小さい加数で埋めても、
\begin{align*}
  x
  &\leq2U_{k+1}+2U_k\\
  &=2(4^{k+1}-2^k)+2(4^k-2^{k-1})\\
  &=10\cdot4^k-3\cdot2^k
  <10\cdot4^k-2=x
\end{align*}
となり、矛盾する。したがって $\lambda\geq3$ である。しかしこのとき、大きい加数3個だけでも、その和は \eqref{eq:three-g} により
\[
  3g_{k+1}=10\cdot4^k-1=x+1
\]
以上となり、再び矛盾する。ゆえに $x\notin F_4(\Pset)$ である。
\end{proof}

\begin{proof}[定理~\ref{thm:asymptotic-order} の証明]
$k\geq2$ を固定し、$m=x_k=10\cdot4^k-2$ とおく。このとき、
\[
  N_k=4m+2=10\cdot4^{k+1}-6
\]
である。補題~\ref{lem:obstruction} により、
\[
  m\notin F_4(\Pset),\qquad m-1\notin F_2(\Pset)
\]
である。したがって、5項判定条件である補題~\ref{lem:fivecrit} から $r(N_k)\geq6$ を得る。整数 $N_k$ は定理~\ref{thm:classification} の11個の例外のいずれでもないので、$r(N_k)\leq6$ である。ゆえに $r(N_k)=6$ である。$N_k\to\infty$ なので、十分大きなすべての偶数 $N$ に対して $r(N)$ を抑える $h\leq5$ は存在しない。一方、定理~\ref{thm:classification} は上界 $h=6$ を与える。したがって $h_{\mathrm{asym}}=6$ である。
\end{proof}

\begin{corollary}[$\Pset$ の位数]\label{cor:orderfive}
集合 $\Pset$ は、非負整数に対する正確に位数 $5$ の漸近加法基底である。より正確には、任意の整数 $n\geq212$ は $F_5(\Pset)$ に属する一方、$F_4(\Pset)$ と $F_3(\Pset)$ はそれぞれ無限に多くの整数を含まない。
\end{corollary}

\begin{proof}
命題~\ref{prop:fiveP} から $[212,\infty)\subseteq F_5(\Pset)$ を得る。任意の $k\geq2$ に対して、補題~\ref{lem:obstruction} は
\[
  x_k=10\cdot4^k-2\notin F_4(\Pset)
\]
を与える。$x_k\to\infty$ なので、漸近的にも4項では不十分である。$F_3(\Pset)\subseteq F_4(\Pset)$ なので、3項についても同様である。
\end{proof}

\begin{corollary}[5項を必要とする4の倍数]
\label{cor:five-multiples}
任意の整数 $k\geq3$ に対し、
\[
  M_k=10\cdot4^{k+1}-8
\]
は $r(M_k)=5$ を満たす。したがって、正確に5個の正の原始Dyck加数を必要とする $4$ の倍数は無限に存在する。
\end{corollary}

\begin{proof}
$M_k=4x_k$ と書く。$k\geq3$ なので $x_k\geq212$ であり、命題~\ref{prop:fiveP} から $x_k\in F_5(\Pset)$ を得る。その表示を $4$ 倍すれば、$r(M_k)\leq5$ である。

$M_k$ が高々4個の原始Dyck加数を用いる表示をもつと仮定し、$2$ の個数を $s$ とする。$M_k\equiv0\pmod4$ なので $s$ は偶数であり、$s\in\{0,2,4\}$ である。$s=0$ なら、$4$ で割ることにより $x_k\in F_4(\Pset)$ を得るが、これは補題~\ref{lem:obstruction} に反する。$s=2$ なら、二つの $2$ を除いてから割ることにより $x_k-1\in F_2(\Pset)$ を得るが、これも矛盾である。$s=4$ なら他の加数を用いる余地がなく、$4$ で割ると $x_k=2$ となるが、これは不可能である。したがって $r(M_k)\geq5$ であり、等号が従う。
\end{proof}

\begin{remark}[二つの合同類]
$N=4n$ の場合、十分大きなすべての $n$ は $F_5(\Pset)$ に属するので、最終的に $r(N)\leq5$ である。系~\ref{cor:five-multiples} は、等号が無限に多く成立することを示す。一方、$N=4m+2$ の場合、$2$ の個数は奇数でなければならない。$m=x_k$ のとき、$2$ を1個用いる場合には $m\in F_4(\Pset)$ が必要であり、3個用いる場合には $m-1\in F_2(\Pset)$ が必要である。補題~\ref{lem:obstruction} はその両方を排除する。この合同類が、相対漸近位数を $5$ から $6$ に引き上げる。
\end{remark}

\section{OEIS数列と計算機検証}
\label{sec:oeis}

三つのOEIS数列が自然に現れる。完全均衡数は \seqnum{A014486}~\cite{OEIS-A014486}、正の原始Dyck数は \seqnum{A057547}~\cite{OEIS-A057547} をなし、本論文の中心数列は
\[
  a(n)=r(2n)
\]
すなわち \seqnum{A395858}~\cite{OEIS-A395858} である。定理~\ref{thm:classification} は6を超えるすべての項を決定する：
\[
  a(n)=8\iff n=23
\]
であり、
\[
  a(n)=7\iff
  n\in\{17,22,49,77,99,101,103,419,421,423\}
\]
である。定理~\ref{thm:asymptotic-order} は、値6をとる明示的な無限部分列
\begin{equation}\label{eq:A395858-six-subsequence}
  a(5\cdot4^{k+1}-3)=6\qquad(k\geq2)
\end{equation}
を与える。

\href{https://github.com/growupkuriyama-hub/lean_cfg_project/blob/3525cfb/LeanCfgProject/PrimDyck/A395858.py}{コミット \texttt{3525cfb} における \path{A395858.py}}
は、表示した初項を再現し、A395858の最初の $50{,}000$ 項を生成する。同じ固定コミットにある補助検証スクリプトは、五つの有限区間包含、命題~\ref{prop:fiveP} の鋭い下端点、および有限例外集合証明書を検査する。その独立な動的計画法による計算は、選択した有限上界まで（既定値では $50{,}000$ まで）の $r(N)$ を検査する。既定の実行では、$2\leq k\leq5$ に対する反復障害、同じ範囲に対する6項族
\[
 r(10\cdot4^{k+1}-6)=6
\]
および $3\leq k\leq5$ に対する5項族
\[
 r(10\cdot4^{k+1}-8)=5
\]
も検査する。この版は、コミット \texttt{3525cfb} におけるリポジトリのPython CI検査を通過した（Python CI run \#1）。これらの計算は再現可能な有限範囲の交差検証を与える。すべての $k$ および十分大きなすべての整数に対する主張は、上で理論的に証明されている。

\section{結語}

本論文の結果は、原始Dyck数に関係する三つの正確な位数を決定する。すべての非負偶数上の相対加法位数は $8$、相対漸近位数は $6$、$\Pset$ の漸近加法位数は $5$ である。長い5項区間を伝播させるのと同じ基数4の自己相似性が、反復障害 $x_k=10\cdot4^k-2$ も生み出す。

$r(N)=6$ となる偶数 $N$ の完全な集合を特徴づけることは、今後の課題として残る。計算は $10\cdot4^k$ の近傍に追加のブロック構造が存在することを示唆するが、それらのブロックの正確な幅と端点はまだ証明されていない。より広い問いとして、他の分解不能なディジタル言語に対しても、正規部分族による上界と世代間の隙間による下界が同様に現れるかが挙げられる。

\begin{thebibliography}{9}

\bibitem{BHS}
J.~P. Bell, K.~G. Hare, and J.~Shallit,
When is an automatic set an additive basis?,
\emph{Proc. Amer. Math. Soc. Ser. B} \textbf{5} (2018), 50--63.

\bibitem{BLS}
J.~P. Bell, T.~F. Lidbetter, and J.~Shallit,
Additive number theory via approximation by regular languages,
\emph{Internat. J. Found. Comput. Sci.} \textbf{31} (2020), 667--687.

\bibitem{MNRS}
P.~Madhusudan, D.~Nowotka, A.~Rajasekaran, and J.~Shallit,
Lagrange's theorem for binary squares,
in \emph{43rd International Symposium on Mathematical Foundations of
Computer Science}, Leibniz Int. Proc. Inform., Vol.~117, Schloss
Dagstuhl--Leibniz-Zentrum f{\"u}r Informatik, 2018, pp.~18:1--18:14.

\bibitem{Nathanson}
M.~B. Nathanson,
\emph{Additive Number Theory: The Classical Bases},
Grad. Texts in Math., Vol.~164, Springer, 1996.

\bibitem{OEIS-A014486}
OEIS Foundation Inc.,
\emph{The On-Line Encyclopedia of Integer Sequences, A014486},
\url{https://oeis.org/A014486}.

\bibitem{OEIS-A057547}
OEIS Foundation Inc.,
\emph{The On-Line Encyclopedia of Integer Sequences, A057547},
\url{https://oeis.org/A057547}.

\bibitem{OEIS-A395858}
OEIS Foundation Inc.,
\emph{The On-Line Encyclopedia of Integer Sequences, A395858},
\url{https://oeis.org/A395858}.

\bibitem{RSS}
A.~Rajasekaran, J.~Shallit, and T.~Smith,
Sums of palindromes: an approach via automata,
in \emph{35th Symposium on Theoretical Aspects of Computer Science},
Leibniz Int. Proc. Inform., Vol.~96, Schloss Dagstuhl--Leibniz-Zentrum
f{\"u}r Informatik, 2018, pp.~54:1--54:12.

\end{thebibliography}

\begin{flushleft}
Takayuki Kuriyama\\
Growup Learning Academy\\
4-8-3 Shakujii-machi\\
Nerima-ku, Tokyo 177-0041\\
Japan\\
\textit{メールアドレス}: \href{mailto:growup.kuriyama@gmail.com}
{\texttt{growup.kuriyama@gmail.com}}
\end{flushleft}

\end{document}
